\section{本章纲要：强国的逻辑与帝国的难题}

从\eventref{WQ-0453-ZHI}{三家分晋}到\eventref{QIN-0209-FALL}{秦亡}，战国秦的长线不是“谁最会打仗”，而是谁最能把社会资源持续转化为国家行动。魏、楚、韩、秦等国的改革共同抬高了竞争门槛；秦把这套能力结合得最彻底，因而完成统一；它又在和平时期以过高强度使用这套能力，因而迅速失去承受者。

下一章的汉，不会简单“推翻秦的一切”。它会继承郡县、法律、道路和大一统的想象，同时寻找一种比秦更能长期维持的统治节奏。
